\section{Relation calculations, splitting functors and determinant limits}
\label{rel:details}
\subsection{Correct graphs and stable spaces}
In the original basis let $f e_1=e_2$, $f e_2=e_3$, $f e_3=0$ and
$g e_1=0$, $g e_2=e_1$, $g e_3=e_2$. Then
\[
\begin{split}
 \Gamma_f&=\langle(e_1,e_2),(e_2,e_3),(e_3,0)\rangle,\\
 \Gamma_{f^2}&=\langle(e_1,e_3),(e_2,0),(e_3,0)\rangle,\\
 \Gamma_{gf^2}&=\langle(e_1,e_2),(e_2,0),(e_3,0)\rangle.
\end{split}
\]
Each graph has dimension three. Omitting zero-output generators changes
the relation into a partial graph and changes its powers and domain.
The following are subspaces of $V$, not lists of pairs in $V\oplus V$:
\[
\begin{array}{c|cc}
T&\bigcup_r\ker T^r&\bigcap_r\operatorname{im}T^r\\\hline
f,g&V&0\\
gf=\operatorname{diag}(1,1,0)&ke_3&ke_1\oplus ke_2\\
fg=\operatorname{diag}(0,1,1)&ke_1&ke_2\oplus ke_3
\end{array}
\]
For all four full graphs $\bigcup_r\Gamma_T^r(0)=0$. This corrects both
the source's graph convention and its interchanged $gf,fg$ kernel entries.
For an arbitrary multivalued relation there need not even be a linear map
on its domain; its regular-part theorem requires the appropriate relation
subquotients and finiteness conditions. We do not replace those hypotheses
by the graph calculation.
For comparison, in \cite[Section 4, Lemmas 4.2,4.5--4.6]{CrawleyBoeveyStrings}
write $C'=\bigcup_n C^n(0)$ and let $C''$ consist of vectors with a
compatible infinite backward chain. In finite dimension $C''=\bigcap_n C^nV$;
this equality is not automatic for arbitrary infinite-dimensional spaces.
The regular quotient is $C^\sharp/C^\flat$, where
\[
 C^\sharp=C''\cap(C^{-1})'',\qquad
 C^\flat=(C''\cap(C^{-1})')+(C'\cap(C^{-1})'').
\]
The automorphism statement uses this quotient for an arbitrary relation.
The splitting statement additionally requires its finite-dimensionality,
rather than the source's unqualified $\iota/\kappa$ expression.

\subsection{Proof of the dessin-category equivalence}
Write $a_i:V_i\to V_{i-1}$ for $i=1,\ldots,m-1$ and
$p_i=a_1\cdots a_i:V_i\to V_0$, with $p_0=1$; the rightmost map acts
first. Every $p_i$ is an isomorphism. Conjugating the arrow maps by these
identifications makes $a_i=1$ for $i\ne0$, and makes
$T=p_{m-1}a_0$ at the remaining arrow. All $b$-relations vanish.
An intertwiner is transported to maps $h_i$ with $h_i=h_{i-1}$ for
$i\ne0$ and $h_{m-1}T=T'h_0$. Hence every $h_i=h$ and $hT=T'h$.
This proves full faithfulness and essential surjectivity, with inverse
given by $(V,T)\mapsto N_m(T)$. It is also an exact equivalence: kernels
and cokernels of intertwiners are the usual invariant kernels/cokernels;
the identity arrow maps remain identities. Finite dimension is a condition
on modules, while $A_m$ itself is infinite dimensional through its
arbitrarily long permitted branch paths.

The sum of the full permitted turns on both branches is central in this
closed dessin algebra. On the chosen modules the $b$-turns vanish and each
$a$-turn equals $T$, so $u$ acts as $T$ at every vertex. If $T^n=0$,
the action factors through $A_m/(u^n)$ and extends to the completion.
If $T$ has a nonzero invertible Fitting summand, no power of $u$ kills
that summand, precluding continuity. For $T=J_r(\lambda)$ with
$\lambda\ne0$, the obstruction is stronger: $u-\lambda$ is a unit of
$k[[u]]$ but its $r$th power annihilates the proposed module, impossible
for a nonzero unital module. This rules out even an unqualified algebraic
extension of that band, not just a chosen topology.

\subsection{Functoriality and the actual filtration}
Kernel dimensions are nondecreasing and at most $n$. Once two successive
kernels agree, applying $T^{-1}$ shows that all later kernels agree. Thus
they stabilize by $n$; rank-nullity gives image stabilization as well.
On $K(T)$, $T^n=0$. On $I(T)$, $T(I(T))=I(T)$; finite dimension makes
this a bijection. If $hT=T'h$, induction gives $hT^r=T'^rh$, hence
\[
 h(\ker T^r)\subseteq\ker T'^r,\qquad
 h(\operatorname{im}T^r)\subseteq\operatorname{im}T'^r.
\]
This proves functoriality even when $\dim V\ne\dim V'$: choose a common
stabilization exponent for the two finite objects. The decomposition
projection onto $K(T)$ along $I(T)$ is canonical and commutes with all
intertwiners; the existence of a Jordan basis is not needed.
Finite dimension is essential. On $V=k[t]$, multiplication by $t$ has
stable kernel zero and stable image $\bigcap_r t^rk[t]=0$, although
$V\ne0$. The same construction $N_m(V,t)$ is an actual infinite-dimensional
module of $A_m$. Thus the finite Fitting split cannot be extended to every
module of our infinite algebra merely by repeating its notation.

Let $k_r=\dim\ker T^r$, with $k_0=0$, and form
$G_r=\ker T^r/\ker T^{r-1}$. For a nilpotent Jordan block of size $s$,
$k_r=\min(r,s)$, so $\dim G_r$ counts the blocks of size at least $r$,
and $2k_r-k_{r-1}-k_{r+1}$ counts blocks of size exactly $r$.
For $J_3(0)\oplus J_1(0)\oplus J_2(1)$ the stable dimensions are
$4$ and $2$; the kernel dimensions begin $0,2,3,4,4$.
The $G_r$ are functorial subquotients. Their dimensions do not specify a
natural choice of basis or split each filtration step.

\begin{proposition}[What the source's splitting-functor assumptions imply]
Let $\mathcal A,\mathcal B_i$ be additive categories with finite
biproducts. Suppose additive functors $S_i:\mathcal B_i\to\mathcal A$
and $F_i:\mathcal A\to\mathcal B_i$ satisfy
$F_iS_i\cong1$, $F_iS_j\cong0$ for $i\ne j$; only finitely many
$F_iM$ are nonzero for each $M$; the $F_i$ jointly reflect isomorphisms;
and maps $\gamma_{i,M}:S_iF_iM\to M$ have invertible $F_i$-images.
Then $M\cong\bigoplus_i S_iF_iM$. Every nonzero indecomposable is
$S_iB$ for a unique index and an indecomposable $B$ up to isomorphism.
\end{proposition}
\begin{proof}
Sum the finitely many nonzero $\gamma_{i,M}$ to a map $\gamma_M$.
After applying $F_j$, all off-diagonal summands vanish, leaving the
assumed isomorphism $F_j\gamma_{j,M}$, or $0\to0$. Joint reflection
therefore makes $\gamma_M$ invertible. For an indecomposable $M$, exactly
one of the nonzero summands can occur. Conversely, if $S_iB=X\oplus Y$
and $B$ is indecomposable, applying $F_i$ makes one summand vanish, say
$F_iX=0$; every other $F_jX$ also vanishes. Joint reflection applied to
$0\to X$ gives $X=0$ up to isomorphism. Finally apply $F_i$ to an
isomorphism $S_iB\cong S_jB'$ to recover both index and object.
\end{proof}
The direction of $\gamma$ matters: it is into $M$, contrary to one source
diagram. Local finiteness is objectwise in $i$, not a finiteness assertion
over all objects of a category. Abelian categories require kernels,
cokernels and the usual exactness axioms; addition is on Hom groups, not
on an alleged group of quiver vertices. The abstract assumptions above
do not assert that the family $F_i$ is faithful on arbitrary morphisms.
Our Fitting example is concrete: $\mathcal B_0$ consists of nilpotent
pairs, $\mathcal B_\times$ of invertible pairs, $S$ are inclusions and
$F_0,F_\times$ the two stable summands.

\subsection{Strings, bands, Hom and one non-rigid string}
In this branch category, Jordan normal form classifies indecomposables
over an algebraically closed field. Over a general field replace Jordan
eigenvalues by powers of monic irreducible polynomials. A Jordan block
$J_r(0)$ is one directed string with $mr$ basis visits and $mr-1$ nonzero
successor arrows; $J_r(\lambda)$ for $\lambda\ne0$ is a permitted
cycle band. Formal inverse letters are reverse traversals, not inverse
elements of $A_m$. Strings and their reverses describe the same string
module under the usual isomorphism convention. Arbitrary infinite modules
are not all classified by this finite calculation; the actual finiteness
hypotheses in \cite{CrawleyBoeveyStrings} remain essential.

Identify a Jordan pair with $k[t]/(t-\lambda)^r$. An intertwiner from
size $r$ to size $s$ is determined by the image of $1$, annihilated by
$(t-\lambda)^r$. For equal eigenvalues this subspace has dimension
$\min(r,s)$. For different eigenvalues, Bezout's identity makes
$t-\lambda$ invertible on the target, so only zero remains. Full
faithfulness above transfers this Hom computation to $A_m$.
The sequence $0\to J_1(0)\to J_2(0)\to J_1(0)\to0$ is not split:
with $J_2(0)e_1=e_2$, a lift of the quotient generator has nonzero first
coordinate and is not killed by $J_2(0)$. Its image under $N_m$ is a
nonsplit ambient-$A_m$ self-extension. Thus even the string
$N_m(J_1(0))$ is not rigid. No assertion that this exact full subcategory
is derived fully faithful is needed for the nonsplitting argument.

For this particular closed $A_m$, every oriented or inverse letter has
exactly two admissible successors: continue on its branch in the same
orientation, or switch branch and reverse orientation. Cancellation and
forbidden pairs, including inverse forbidden pairs, remove the other
choices. A fixed ordering of these choices identifies each first-letter
cylinder with $\{0,1\}^{\mathbb N}$. The whole infinite-word space is a
finite disjoint union of Cantor spaces, hence itself Cantor. There are no
terminating words inside that space to remove. Arbitrary gentle quivers
can have isolated or eventually forced paths, so this conclusion is not
universal. Algebraic covers retain the lifted zero-relations and deck
equivariance data; arbitrary periodic cover representations are not
automatically bands or a unitary time evolution.

\subsection{A determinant comparison with its exact limits}
On $\bigoplus_iV_i$, $L^m=\operatorname{diag}(T,\ldots,T)$.
For a direct determinant proof of~\eqref{rel:det}, successively eliminate
the blocks $1$ at vertices $1,\ldots,m-1$ in $1-zL$.
Block row additions preserve determinant and leave $1-z^mT$ at vertex
zero. This works over $k[z]$ without inverting $z$ or $T$; the zero-size
case gives one on both sides. The Fitting splitting separates a nilpotent
factor $1$ and the determinant on the invertible summand. For instance
$m=5$, $T=J_2(2)$ gives $(1-2z^5)^2$, whose complex roots have modulus
$2^{-1/5}$, despite being a genuine genus-two dessin module.

Under conjugation of $T$ this is an invariant polynomial. With independent
source/target bases a presentation determinant is instead a semi-invariant
with its determinant-character weight; those actions must not be conflated.
The polynomial equality does not construct a positive Hermitian metric,
an arithmetic Euler factor, a regularized determinant or a zero-preserving
map to completed $L$-functions. Functorial kernel filtrations similarly do
not construct an equivalence of derived/Fukaya categories. A fixed graded
surface category and a determinant-compatible arithmetic realization remain
separate obligations.

The exact rational checker evaluates actual genus-two/four quivers, Hom
equations, full-turn matrices, both filtration directions and the nonsplit
extension. \path{lean-analytic/LinearRelations.lean} proves full graph
identities, semiconjugacy and filtration preservation, actual finite-space
Fitting splitting/stabilization and unit-power annihilation. It does not
package the entire dessin category, Jordan classification, general
splitting-functor proposition or determinant proof in Lean.
